\documentclass{beamer}
\usepackage[utf8]{inputenc}
\usepackage[spanish]{babel}
\usepackage{listings}
\usepackage{xcolor}
\usepackage{tikz}
\usepackage{tcolorbox}
\usepackage{fontawesome5}

% Configuración de colores - Gruvbox Light con marrón principal
\definecolor{bg0}{HTML}{fbf1c7}
\definecolor{bg1}{HTML}{ebdbb2}
\definecolor{bg2}{HTML}{d5c4a1}
\definecolor{bg3}{HTML}{bdae93}
\definecolor{fg0}{HTML}{282828}
\definecolor{fg1}{HTML}{3c3836}
\definecolor{fg2}{HTML}{504945}
\definecolor{aqua}{HTML}{689d6a}
\definecolor{aquabright}{HTML}{427b58}
\definecolor{blue}{HTML}{458588}
\definecolor{red}{HTML}{cc241d}
\definecolor{orange}{HTML}{d65d0e}
\definecolor{yellow}{HTML}{d79921}
\definecolor{green}{HTML}{98971a}
\definecolor{purple}{HTML}{b16286}
\definecolor{brown}{HTML}{AB886D}
\definecolor{brownbright}{HTML}{8B6849}

% Colores para código
\definecolor{codegreen}{HTML}{98971a}
\definecolor{codegray}{HTML}{7c6f64}
\definecolor{codepurple}{HTML}{b16286}
\definecolor{backcolour}{HTML}{f9f5d7}

% Configuración de código
\lstdefinestyle{pythonstyle}{
    backgroundcolor=\color{backcolour},   
    commentstyle=\color{codegreen},
    keywordstyle=\color{brownbright}\bfseries,
    numberstyle=\tiny\color{codegray},
    stringstyle=\color{codepurple},
    basicstyle=\ttfamily\scriptsize\color{fg0},
    breakatwhitespace=false,         
    breaklines=true,                 
    captionpos=b,                    
    keepspaces=true,                 
    numbers=left,                    
    numbersep=2pt,                  
    showspaces=false,                
    showstringspaces=false,
    showtabs=false,                  
    tabsize=2,
    language=Python,
    frame=single,
    rulecolor=\color{brown},
    xleftmargin=6pt,
    xrightmargin=6pt,
    inputencoding=utf8,
    extendedchars=true,
    literate=
        {á}{{\'a}}1 {é}{{\'e}}1 {í}{{\'i}}1 {ó}{{\'o}}1 {ú}{{\'u}}1
        {Á}{{\'A}}1 {É}{{\'E}}1 {Í}{{\'I}}1 {Ó}{{\'O}}1 {Ú}{{\'U}}1
        {ñ}{{\~n}}1 {Ñ}{{\~N}}1
        {¿}{{?`}}1 {¡}{{!`}}1
}

\lstset{style=pythonstyle}

% Tema personalizado con marrón principal
\usetheme{Madrid}
\setbeamercolor{structure}{fg=brown}
\setbeamercolor{palette primary}{bg=brown,fg=bg0}
\setbeamercolor{palette secondary}{bg=brownbright,fg=bg0}
\setbeamercolor{palette tertiary}{bg=fg1,fg=bg0}
\setbeamercolor{palette quaternary}{bg=brown,fg=bg0}
\setbeamercolor{block title}{bg=brown,fg=bg0}
\setbeamercolor{block body}{bg=bg1,fg=fg0}
\setbeamercolor{frametitle}{bg=brown,fg=bg0}
\setbeamercolor{title}{fg=fg0}
\setbeamercolor{subtitle}{fg=fg1}
\setbeamercolor{normal text}{fg=fg0,bg=bg0}
\setbeamercolor{item}{fg=brown}
\setbeamercolor{section in toc}{fg=brown}

% Cajas personalizadas con marrón
\newtcolorbox{conceptbox}[1]{
    colback=bg1,
    colframe=brown,
    fonttitle=\bfseries\small,
    title=#1,
    coltitle=fg0,
    left=3pt,
    right=3pt,
    top=3pt,
    bottom=3pt,
    boxsep=2pt
}

\newtcolorbox{notebox}{
    colback=yellow!15,
    colframe=orange,
    leftrule=2mm,
    rightrule=0mm,
    toprule=0mm,
    bottomrule=0mm,
    arc=0mm,
    left=3pt,
    right=3pt,
    top=2pt,
    bottom=2pt
}

\title{\textbf{Introducción a la Programación con Python}}
\subtitle{Semana 8: Programación Orientada a Objetos}
\author{$\nabla$Nablabs}
\date{}
\begin{document}

\begin{frame}
\titlepage
\end{frame}

\begin{frame}
\frametitle{Contenido}
\tableofcontents
\end{frame}

\section{¿Qué es la POO?}

\begin{frame}{Programación Orientada a Objetos \faCube}
\begin{conceptbox}{¿Qué es POO?}
\small
\begin{itemize}
    \item[\faBox] Paradigma de programación basado en \textbf{objetos}
    \item[\faPuzzlePiece] Objetos combinan \textbf{datos} (atributos) y \textbf{comportamiento} (métodos)
    \item[\faLayerGroup] Organiza código de forma modular y reutilizable
    \item[\faProjectDiagram] Modela conceptos del mundo real
    \item[\faShieldAlt] Encapsula datos y funcionalidad
\end{itemize}
\end{conceptbox}

\vspace{0.2cm}
\begin{center}
\large\textcolor{brown}{\faLightbulb\ ¡Piensa en objetos, no solo en funciones!}
\end{center}
\end{frame}

\begin{frame}{¿Por qué usar POO? \faQuestion}
\begin{columns}[T]
\column{0.48\textwidth}
\textbf{Sin POO:} \faTimesCircle
\small
\begin{itemize}
    \item Datos dispersos
    \item Funciones aisladas
    \item Difícil de mantener
    \item Código repetitivo
    \item Menos organización
\end{itemize}

\column{0.48\textwidth}
\textbf{Con POO:} \faCheckCircle
\small
\begin{itemize}
    \item Datos agrupados
    \item Comportamiento asociado
    \item Fácil de extender
    \item Código reutilizable
    \item Mejor organización
\end{itemize}
\end{columns}

\vspace{0.3cm}
\begin{center}
\textcolor{brown}{\faInfoCircle\ POO facilita modelar sistemas complejos}
\end{center}
\end{frame}

\section{Clases y objetos}

\begin{frame}[fragile]{Clases: El molde \faDrawPolygon}
\begin{conceptbox}{¿Qué es una clase?}
Una \textbf{clase} es un molde o plantilla para crear objetos. Define qué atributos y métodos tendrán los objetos.
\end{conceptbox}

\textbf{Ejemplo conceptual:}
\small
\begin{itemize}
    \item[\faDrawPolygon] \textbf{Clase}: Perro (el concepto general)
    \item[\faDog] \textbf{Objeto}: Mi perro "Max" (una instancia específica)
\end{itemize}

\vspace{0.2cm}
\textbf{Sintaxis básica:}
\begin{lstlisting}
class NombreClase:
    # Definición de la clase
    pass
\end{lstlisting}

\vspace{-0.1cm}
\begin{notebox}
\small\faInfoCircle\ Por convención, los nombres de clases usan CamelCase
\end{notebox}
\end{frame}

\begin{frame}[fragile]{Primera clase simple \faCode}
\begin{lstlisting}
class Perro:
    # Atributos de clase
    especie = "Canis familiaris"
    
    # Método simple
    def ladrar(self):
        return "¡Guau guau!"

# Crear objetos (instancias)
mi_perro = Perro()
tu_perro = Perro()

print(mi_perro.especie)
print(mi_perro.ladrar())
\end{lstlisting}

\vspace{-0.2cm}
\begin{tcolorbox}[colback=green!10,colframe=green,title={\small\faArrowCircleRight \ Salida},coltitle=fg0,left=3pt,right=3pt,top=2pt,bottom=2pt,boxsep=1pt]
\small\texttt{Canis familiaris\\
¡Guau guau!}
\end{tcolorbox}
\end{frame}

\begin{frame}[fragile]{El parámetro self \faUser}
\begin{conceptbox}{¿Qué es self?}
\textbf{self} es una referencia a la instancia actual del objeto. Permite acceder a sus atributos y métodos.
\end{conceptbox}

\begin{lstlisting}
class Persona:
    def saludar(self):
        return "Hola, soy una persona"
    
    def presentarse(self, nombre):
        # self permite acceder al mismo objeto
        return f"{self.saludar()}, me llamo {nombre}"

persona = Persona()
print(persona.presentarse("Ana"))
\end{lstlisting}

\vspace{-0.2cm}
\begin{tcolorbox}[colback=green!10,colframe=green,title={\small\faArrowCircleRight \ Salida},coltitle=fg0,left=3pt,right=3pt,top=2pt,bottom=2pt,boxsep=1pt]
\small\texttt{Hola, soy una persona, me llamo Ana}
\end{tcolorbox}

\vspace{-0.1cm}
\begin{notebox}
\small\faExclamationTriangle\ self debe ser el primer parámetro de los métodos de instancia
\end{notebox}
\end{frame}

\section{Método \_\_init\_\_ (constructor)}

\begin{frame}[fragile]{Constructor: \_\_init\_\_ \faWrench}
\begin{conceptbox}{¿Qué es \_\_init\_\_?}
El \textbf{constructor} es un método especial que se ejecuta automáticamente cuando se crea un objeto. Inicializa los atributos.
\end{conceptbox}

\begin{lstlisting}
class Persona:
    def __init__(self, nombre, edad):
        self.nombre = nombre  # Atributo de instancia
        self.edad = edad      # Atributo de instancia
    
    def presentarse(self):
        return f"Hola, soy {self.nombre} y tengo {self.edad} años"

# Crear objetos con valores iniciales
ana = Persona("Ana", 25)
luis = Persona("Luis", 30)

print(ana.presentarse())
print(luis.presentarse())
\end{lstlisting}
\end{frame}

\begin{frame}[fragile]{Constructor: Salida \faArrowCircleRight}
\begin{tcolorbox}[colback=green!10,colframe=green,title={\small\faArrowCircleRight \ Salida},coltitle=fg0,left=3pt,right=3pt,top=3pt,bottom=3pt,boxsep=2pt]
\small\texttt{Hola, soy Ana y tengo 25 años\\
Hola, soy Luis y tengo 30 años}
\end{tcolorbox}

\vspace{0.3cm}
\begin{notebox}
\small\faLightbulb\ Cada objeto tiene sus propios valores de atributos
\end{notebox}
\end{frame}

\section{Atributos de instancia y de clase}

\begin{frame}[fragile]{Atributos de instancia \faUser}
\begin{conceptbox}{¿Qué son?}
Atributos \textbf{únicos} para cada objeto. Se definen en \_\_init\_\_ con self.
\end{conceptbox}

\begin{lstlisting}
class Coche:
    def __init__(self, marca, modelo, año):
        # Atributos de instancia
        self.marca = marca
        self.modelo = modelo
        self.año = año
        self.kilometraje = 0  # Valor inicial por defecto

coche1 = Coche("Toyota", "Corolla", 2020)
coche2 = Coche("Honda", "Civic", 2021)

print(f"{coche1.marca} {coche1.modelo}")
print(f"{coche2.marca} {coche2.modelo}")
\end{lstlisting}

\vspace{-0.2cm}
\begin{tcolorbox}[colback=green!10,colframe=green,title={\small\faArrowCircleRight \ Salida},coltitle=fg0,left=3pt,right=3pt,top=2pt,bottom=2pt,boxsep=1pt]
\small\texttt{Toyota Corolla\\
Honda Civic}
\end{tcolorbox}
\end{frame}

\begin{frame}[fragile]{Atributos de clase \faLayerGroup}
\begin{conceptbox}{¿Qué son?}
Atributos \textbf{compartidos} por todas las instancias de la clase. Se definen fuera de \_\_init\_\_.
\end{conceptbox}

\begin{lstlisting}
class Empleado:
    # Atributo de clase (compartido)
    empresa = "TechCorp"
    total_empleados = 0
    
    def __init__(self, nombre, salario):
        # Atributos de instancia (únicos)
        self.nombre = nombre
        self.salario = salario
        Empleado.total_empleados += 1

emp1 = Empleado("Ana", 50000)
emp2 = Empleado("Luis", 60000)

print(f"Empresa: {Empleado.empresa}")
print(f"Total empleados: {Empleado.total_empleados}")
\end{lstlisting}
\end{frame}

\begin{frame}[fragile]{Atributos de clase: Salida \faArrowCircleRight}
\begin{tcolorbox}[colback=green!10,colframe=green,title={\small\faArrowCircleRight \ Salida},coltitle=fg0,left=3pt,right=3pt,top=3pt,bottom=3pt,boxsep=2pt]
\small\texttt{Empresa: TechCorp\\
Total empleados: 2}
\end{tcolorbox}

\vspace{0.3cm}
\begin{notebox}
\small\faInfoCircle\ Los atributos de clase se acceden con NombreClase.atributo
\end{notebox}
\end{frame}

\section{Métodos de instancia}

\begin{frame}[fragile]{Métodos de instancia \faTools}
\begin{conceptbox}{¿Qué son?}
Funciones definidas dentro de una clase que operan sobre los atributos de la instancia. Siempre llevan \textbf{self} como primer parámetro.
\end{conceptbox}

\begin{lstlisting}
class CuentaBancaria:
    def __init__(self, titular, saldo=0):
        self.titular = titular
        self.saldo = saldo
    
    def depositar(self, monto):
        self.saldo += monto
        return f"Depósito exitoso. Saldo: ${self.saldo}"
    
    def retirar(self, monto):
        if monto <= self.saldo:
            self.saldo -= monto
            return f"Retiro exitoso. Saldo: ${self.saldo}"
        return "Fondos insuficientes"
\end{lstlisting}
\end{frame}

\begin{frame}[fragile]{Métodos de instancia: Uso \faPlay}
\begin{lstlisting}
# Crear cuenta
cuenta = CuentaBancaria("Ana García", 1000)

print(cuenta.depositar(500))
print(cuenta.retirar(300))
print(cuenta.retirar(2000))
print(f"Saldo final: ${cuenta.saldo}")
\end{lstlisting}

\vspace{-0.2cm}
\begin{tcolorbox}[colback=green!10,colframe=green,title={\small\faArrowCircleRight \ Salida},coltitle=fg0,left=3pt,right=3pt,top=3pt,bottom=3pt,boxsep=2pt]
\small\texttt{Depósito exitoso. Saldo: \$1500\\
Retiro exitoso. Saldo: \$1200\\
Fondos insuficientes\\
Saldo final: \$1200}
\end{tcolorbox}
\end{frame}

\section{Encapsulación y convenciones}

\begin{frame}[fragile]{Encapsulación \faLock}
\begin{conceptbox}{¿Qué es encapsulación?}
Ocultar los detalles internos de un objeto y exponer solo lo necesario. Protege los datos de modificaciones no deseadas.
\end{conceptbox}

\textbf{Convenciones en Python:}
\small
\begin{itemize}
    \item[\faUnlock] \texttt{nombre}: Público (accesible desde cualquier lugar)
    \item[\faKey] \texttt{\_nombre}: Protegido (convención, uso interno)
    \item[\faLock] \texttt{\_\_nombre}: Privado (name mangling)
\end{itemize}

\vspace{0.2cm}
\begin{notebox}
\small\faInfoCircle\ Python no tiene verdadera privacidad, solo convenciones
\end{notebox}
\end{frame}

\begin{frame}[fragile]{Atributos privados \faUserSecret}
\begin{lstlisting}
class CuentaSegura:
    def __init__(self, titular, saldo):
        self.titular = titular        # Público
        self._password = "1234"       # Protegido
        self.__saldo = saldo          # Privado
    
    def consultar_saldo(self, password):
        if password == self._password:
            return f"Saldo: ${self.__saldo}"
        return "Contraseña incorrecta"
    
    def depositar(self, monto):
        self.__saldo += monto

cuenta = CuentaSegura("Ana", 1000)
print(cuenta.consultar_saldo("1234"))
# print(cuenta.__saldo)  # ¡Error! Atributo privado
\end{lstlisting}
\end{frame}

\begin{frame}[fragile]{Atributos privados: Salida \faArrowCircleRight}
\begin{tcolorbox}[colback=green!10,colframe=green,title={\small\faArrowCircleRight \ Salida},coltitle=fg0,left=3pt,right=3pt,top=3pt,bottom=3pt,boxsep=2pt]
\small\texttt{Saldo: \$1000}
\end{tcolorbox}

\vspace{0.3cm}
\begin{notebox}
\small\faExclamationTriangle\ Los atributos privados usan "name mangling": \_NombreClase\_\_atributo
\end{notebox}
\end{frame}

\section{Propiedades con @property}

\begin{frame}[fragile]{Decorador @property \faGem}
\begin{conceptbox}{¿Qué es @property?}
Permite crear atributos que se comportan como métodos. Controla el acceso y modificación de atributos.
\end{conceptbox}

\begin{lstlisting}
class Temperatura:
    def __init__(self, celsius):
        self._celsius = celsius
    
    @property
    def celsius(self):
        return self._celsius
    
    @property
    def fahrenheit(self):
        return (self._celsius * 9/5) + 32
    
    @celsius.setter
    def celsius(self, valor):
        if valor < -273.15:
            raise ValueError("Temperatura bajo cero absoluto")
        self._celsius = valor
\end{lstlisting}
\end{frame}

\begin{frame}[fragile]{@property: Uso \faThermometer}
\begin{lstlisting}
temp = Temperatura(25)

# Acceder como atributo (no método)
print(f"Celsius: {temp.celsius}°C")
print(f"Fahrenheit: {temp.fahrenheit}°F")

# Modificar con validación
temp.celsius = 30
print(f"Nueva temp: {temp.celsius}°C")

# temp.celsius = -300  # ¡Error! ValueError
\end{lstlisting}

\vspace{-0.2cm}
\begin{tcolorbox}[colback=green!10,colframe=green,title={\small\faArrowCircleRight \ Salida},coltitle=fg0,left=3pt,right=3pt,top=3pt,bottom=3pt,boxsep=2pt]
\small\texttt{Celsius: 25°C\\
Fahrenheit: 77.0°F\\
Nueva temp: 30°C}
\end{tcolorbox}
\end{frame}

\section{Métodos especiales}

\begin{frame}[fragile]{Métodos especiales (dunder) \faMagic}
\begin{conceptbox}{¿Qué son?}
Métodos con doble guion bajo (\_\_método\_\_) que Python llama automáticamente en situaciones específicas.
\end{conceptbox}

\textbf{Métodos comunes:}
\small
\begin{itemize}
    \item[\faWrench] \texttt{\_\_init\_\_}: Constructor
    \item[\faPrint] \texttt{\_\_str\_\_}: Representación legible (para usuarios)
    \item[\faCode] \texttt{\_\_repr\_\_}: Representación técnica (para desarrolladores)
    \item[\faEquals] \texttt{\_\_eq\_\_}: Igualdad (==)
    \item[\faCalculator] \texttt{\_\_add\_\_}: Suma (+)
    \item[\faListOl] \texttt{\_\_len\_\_}: Longitud (len())
\end{itemize}
\end{frame}

\begin{frame}[fragile]{\_\_str\_\_ y \_\_repr\_\_ \faPrint}
\begin{lstlisting}
class Libro:
    def __init__(self, titulo, autor, año):
        self.titulo = titulo
        self.autor = autor
        self.año = año
    
    def __str__(self):
        # Para usuarios (print)
        return f'"{self.titulo}" por {self.autor}'
    
    def __repr__(self):
        # Para desarrolladores (consola)
        return f"Libro('{self.titulo}', '{self.autor}', {self.año})"

libro = Libro("Cien años de soledad", "García Márquez", 1967)

print(libro)        # Llama __str__
print(repr(libro))  # Llama __repr__
\end{lstlisting}
\end{frame}

\begin{frame}[fragile]{\_\_str\_\_ y \_\_repr\_\_: Salida \faArrowCircleRight}
\begin{tcolorbox}[colback=green!10,colframe=green,title={\small\faArrowCircleRight \ Salida},coltitle=fg0,left=3pt,right=3pt,top=3pt,bottom=3pt,boxsep=2pt]
\small\texttt{"Cien años de soledad" por García Márquez\\
Libro('Cien años de soledad', 'García Márquez', 1967)}
\end{tcolorbox}

\vspace{0.3cm}
\begin{notebox}
\small\faLightbulb\ \_\_str\_\_ para usuarios, \_\_repr\_\_ para debug
\end{notebox}
\end{frame}

\begin{frame}[fragile]{Más métodos especiales \faCalculator}
\begin{lstlisting}
class Punto:
    def __init__(self, x, y):
        self.x = x
        self.y = y
    
    def __str__(self):
        return f"({self.x}, {self.y})"
    
    def __add__(self, otro):
        return Punto(self.x + otro.x, self.y + otro.y)
    
    def __eq__(self, otro):
        return self.x == otro.x and self.y == otro.y

p1 = Punto(1, 2)
p2 = Punto(3, 4)
p3 = p1 + p2  # Llama __add__

print(f"p1 + p2 = {p3}")
print(f"p1 == p2: {p1 == p2}")  # Llama __eq__
\end{lstlisting}
\end{frame}

\begin{frame}[fragile]{Métodos especiales: Salida \faArrowCircleRight}
\begin{tcolorbox}[colback=green!10,colframe=green,title={\small\faArrowCircleRight \ Salida},coltitle=fg0,left=3pt,right=3pt,top=3pt,bottom=3pt,boxsep=2pt]
\small\texttt{p1 + p2 = (4, 6)\\
p1 == p2: False}
\end{tcolorbox}

\vspace{0.3cm}
\begin{notebox}
\small\faInfoCircle\ Los métodos especiales hacen que tus objetos se comporten como tipos nativos
\end{notebox}
\end{frame}

\section{Herencia}

\begin{frame}[fragile]{Herencia \faProjectDiagram}
\begin{conceptbox}{¿Qué es herencia?}
Mecanismo que permite crear clases nuevas basadas en clases existentes, heredando sus atributos y métodos.
\end{conceptbox}

\begin{lstlisting}
class Animal:
    def __init__(self, nombre):
        self.nombre = nombre
    
    def hacer_sonido(self):
        return "Hace un sonido"
    
    def dormir(self):
        return f"{self.nombre} está durmiendo"

# Perro hereda de Animal
class Perro(Animal):
    def hacer_sonido(self):
        return "¡Guau guau!"

# Gato hereda de Animal
class Gato(Animal):
    def hacer_sonido(self):
        return "¡Miau miau!"
\end{lstlisting}
\end{frame}

\begin{frame}[fragile]{Herencia: Uso \faDog}
\begin{lstlisting}
perro = Perro("Max")
gato = Gato("Luna")

print(perro.hacer_sonido())  # Método sobrescrito
print(gato.hacer_sonido())   # Método sobrescrito

print(perro.dormir())  # Método heredado
print(gato.dormir())   # Método heredado
\end{lstlisting}

\vspace{-0.2cm}
\begin{tcolorbox}[colback=green!10,colframe=green,title={\small\faArrowCircleRight \ Salida},coltitle=fg0,left=3pt,right=3pt,top=3pt,bottom=3pt,boxsep=2pt]
\small\texttt{¡Guau guau!\\
¡Miau miau!\\
Max está durmiendo\\
Luna está durmiendo}
\end{tcolorbox}
\end{frame}

\begin{frame}[fragile]{super(): Llamar a la clase padre \faArrowUp}
\begin{lstlisting}
class Vehiculo:
    def __init__(self, marca, modelo):
        self.marca = marca
        self.modelo = modelo
    
    def descripcion(self):
        return f"{self.marca} {self.modelo}"

class Coche(Vehiculo):
    def __init__(self, marca, modelo, puertas):
        super().__init__(marca, modelo)  # Llama al padre
        self.puertas = puertas
    
    def descripcion(self):
        # Reutilizar método del padre y extenderlo
        return f"{super().descripcion()} - {self.puertas} puertas"

coche = Coche("Toyota", "Corolla", 4)
print(coche.descripcion())
\end{lstlisting}
\end{frame}

\begin{frame}[fragile]{super(): Salida \faArrowCircleRight}
\begin{tcolorbox}[colback=green!10,colframe=green,title={\small\faArrowCircleRight \ Salida},coltitle=fg0,left=3pt,right=3pt,top=3pt,bottom=3pt,boxsep=2pt]
\small\texttt{Toyota Corolla - 4 puertas}
\end{tcolorbox}

\vspace{0.3cm}
\begin{notebox}
\small\faLightbulb\ super() evita repetir código y facilita el mantenimiento
\end{notebox}
\end{frame}

\section{Métodos de clase}

\begin{frame}[fragile]{@classmethod \faLayerGroup}
\begin{conceptbox}{¿Qué es @classmethod?}
Método que recibe la clase (cls) como primer argumento, no la instancia. Puede acceder a atributos de clase.
\end{conceptbox}

\begin{lstlisting}
class Producto:
    iva = 0.16  # Atributo de clase
    
    def __init__(self, nombre, precio_base):
        self.nombre = nombre
        self.precio_base = precio_base
    
    @classmethod
    def cambiar_iva(cls, nuevo_iva):
        cls.iva = nuevo_iva
    
    def precio_final(self):
        return self.precio_base * (1 + Producto.iva)

# Cambiar IVA para todos los productos
Producto.cambiar_iva(0.18)
\end{lstlisting}
\end{frame}

\begin{frame}[fragile]{@classmethod: Uso \faTag}
\begin{lstlisting}
# Crear productos después del cambio de IVA
p1 = Producto("Laptop", 1000)
p2 = Producto("Mouse", 20)

print(f"{p1.nombre}: ${p1.precio_final():.2f}")
print(f"{p2.nombre}: ${p2.precio_final():.2f}")

# Cambiar IVA nuevamente
Producto.cambiar_iva(0.20)
print(f"\nNuevo precio {p1.nombre}: ${p1.precio_final():.2f}")
\end{lstlisting}

\vspace{-0.2cm}
\begin{tcolorbox}[colback=green!10,colframe=green,title={\small\faArrowCircleRight \ Salida},coltitle=fg0,left=3pt,right=3pt,top=3pt,bottom=3pt,boxsep=2pt]
\small\texttt{Laptop: \$1180.00\\
Mouse: \$23.60\\
\ \\
Nuevo precio Laptop: \$1200.00}
\end{tcolorbox}
\end{frame}

\begin{frame}[fragile]{@classmethod como constructor alternativo \faWrench}
\begin{lstlisting}
class Fecha:
    def __init__(self, dia, mes, año):
        self.dia = dia
        self.mes = mes
        self.año = año
    
    @classmethod
    def desde_string(cls, fecha_str):
        # Constructor alternativo desde string
        dia, mes, año = map(int, fecha_str.split('/'))
        return cls(dia, mes, año)
    
    def __str__(self):
        return f"{self.dia:02d}/{self.mes:02d}/{self.año}"

fecha1 = Fecha(15, 10, 2025)
fecha2 = Fecha.desde_string("25/12/2025")

print(fecha1)
print(fecha2)
\end{lstlisting}
\end{frame}

\begin{frame}[fragile]{@classmethod alternativo: Salida \faArrowCircleRight}
\begin{tcolorbox}[colback=green!10,colframe=green,title={\small\faArrowCircleRight \ Salida},coltitle=fg0,left=3pt,right=3pt,top=3pt,bottom=3pt,boxsep=2pt]
\small\texttt{15/10/2025\\
25/12/2025}
\end{tcolorbox}

\vspace{0.3cm}
\begin{notebox}
\small\faLightbulb\ @classmethod es útil para constructores alternativos
\end{notebox}
\end{frame}

\section{Métodos estáticos}

\begin{frame}[fragile]{@staticmethod \faTools}
\begin{conceptbox}{¿Qué es @staticmethod?}
Método que \textbf{no} recibe ni self ni cls. Es una función regular agrupada en la clase por organización.
\end{conceptbox}

\begin{lstlisting}
class Matematicas:
    @staticmethod
    def es_par(numero):
        return numero % 2 == 0
    
    @staticmethod
    def factorial(n):
        if n <= 1:
            return 1
        return n * Matematicas.factorial(n - 1)
    
    @staticmethod
    def es_primo(n):
        if n < 2:
            return False
        for i in range(2, int(n ** 0.5) + 1):
            if n % i == 0:
                return False
        return True
\end{lstlisting}
\end{frame}

\begin{frame}[fragile]{@staticmethod: Uso \faCalculator}
\begin{lstlisting}
# No necesitan instancia de la clase
print(f"¿4 es par? {Matematicas.es_par(4)}")
print(f"¿7 es par? {Matematicas.es_par(7)}")

print(f"Factorial de 5: {Matematicas.factorial(5)}")

print(f"¿17 es primo? {Matematicas.es_primo(17)}")
print(f"¿18 es primo? {Matematicas.es_primo(18)}")
\end{lstlisting}

\vspace{-0.2cm}
\begin{tcolorbox}[colback=green!10,colframe=green,title={\small\faArrowCircleRight \ Salida},coltitle=fg0,left=3pt,right=3pt,top=3pt,bottom=3pt,boxsep=2pt]
\small\texttt{¿4 es par? True\\
¿7 es par? False\\
Factorial de 5: 120\\
¿17 es primo? True\\
¿18 es primo? False}
\end{tcolorbox}
\end{frame}

\begin{frame}{Comparación: Métodos \faTable}
\begin{center}
\small
\begin{tabular}{|l|c|c|c|}
\hline
\textbf{Tipo} & \textbf{Primer arg} & \textbf{Acceso a} & \textbf{Uso} \\
\hline
Instancia & \texttt{self} & Instancia & Operaciones del objeto \\
\hline
Clase & \texttt{cls} & Clase & Config/Constructores \\
\hline
Estático & - & Nada & Utilidades agrupadas \\
\hline
\end{tabular}
\end{center}

\vspace{0.3cm}
\begin{notebox}
\small\faInfoCircle\ Usa el tipo de método según lo que necesites acceder
\end{notebox}
\end{frame}

\section{Proyecto: Sistema de gestión}

\begin{frame}{Proyecto: Sistema de gestión de biblioteca \faBook}
\begin{conceptbox}{Objetivo}
Crear un sistema completo usando POO para gestionar una biblioteca con libros, usuarios y préstamos.
\end{conceptbox}

\textbf{Funcionalidades:}
\small
\begin{itemize}
    \item[\faBook] Gestión de libros (agregar, buscar, listar)
    \item[\faUser] Gestión de usuarios
    \item[\faExchangeAlt] Sistema de préstamos y devoluciones
    \item[\faHistory] Historial de préstamos
    \item[\faCheckCircle] Validaciones y restricciones
\end{itemize}

\vspace{0.2cm}
\begin{center}
\textcolor{brown}{\faLightbulb\ Aplicaremos todos los conceptos de POO}
\end{center}
\end{frame}

\begin{frame}[fragile]{Proyecto: Clase Libro \faBook}
\begin{lstlisting}
class Libro:
    total_libros = 0  # Atributo de clase
    
    def __init__(self, titulo, autor, isbn, copias=1):
        self.titulo = titulo
        self.autor = autor
        self.isbn = isbn
        self.copias_totales = copias
        self.copias_disponibles = copias
        Libro.total_libros += 1
    
    def __str__(self):
        return f'"{self.titulo}" por {self.autor}'
    
    def __repr__(self):
        return f"Libro('{self.titulo}', '{self.autor}', '{self.isbn}')"
    
    def esta_disponible(self):
        return self.copias_disponibles > 0
\end{lstlisting}
\end{frame}

\begin{frame}[fragile]{Proyecto: Clase Usuario \faUser}
\begin{lstlisting}
class Usuario:
    def __init__(self, nombre, id_usuario):
        self.nombre = nombre
        self.id_usuario = id_usuario
        self._libros_prestados = []  # Protegido
        self.max_prestamos = 3
    
    def __str__(self):
        return f"{self.nombre} (ID: {self.id_usuario})"
    
    def puede_prestar(self):
        return len(self._libros_prestados) < self.max_prestamos
    
    @property
    def cantidad_prestamos(self):
        return len(self._libros_prestados)
    
    def agregar_prestamo(self, libro):
        self._libros_prestados.append(libro)
    
    def devolver_libro(self, libro):
        self._libros_prestados.remove(libro)
\end{lstlisting}
\end{frame}

\begin{frame}[fragile]{Proyecto: Clase Prestamo \faExchangeAlt}
\begin{lstlisting}
from datetime import datetime, timedelta

class Prestamo:
    def __init__(self, libro, usuario, dias=14):
        self.libro = libro
        self.usuario = usuario
        self.fecha_prestamo = datetime.now()
        self.fecha_devolucion = self.fecha_prestamo + timedelta(days=dias)
        self.devuelto = False
    
    def __str__(self):
        estado = "Devuelto" if self.devuelto else "Activo"
        return f"{self.libro} - {self.usuario} [{estado}]"
    
    def esta_vencido(self):
        if self.devuelto:
            return False
        return datetime.now() > self.fecha_devolucion
    
    def marcar_devuelto(self):
        self.devuelto = True
\end{lstlisting}
\end{frame}

\begin{frame}[fragile]{Proyecto: Clase Biblioteca \faBuilding}
\begin{lstlisting}
class Biblioteca:
    def __init__(self, nombre):
        self.nombre = nombre
        self._catalogo = []      # Libros
        self._usuarios = []       # Usuarios registrados
        self._prestamos = []      # Historial de préstamos
    
    def agregar_libro(self, libro):
        self._catalogo.append(libro)
        print(f"✓ Libro agregado: {libro}")
    
    def registrar_usuario(self, usuario):
        self._usuarios.append(usuario)
        print(f"✓ Usuario registrado: {usuario}")
    
    def buscar_libro(self, titulo):
        for libro in self._catalogo:
            if titulo.lower() in libro.titulo.lower():
                return libro
        return None
\end{lstlisting}

\small(Continuación en siguiente slide...)
\end{frame}

\begin{frame}[fragile]{Proyecto: Biblioteca - Préstamo \faHandHoldingHeart}
\begin{lstlisting}
    def prestar_libro(self, titulo, id_usuario):
        # Buscar libro
        libro = self.buscar_libro(titulo)
        if not libro:
            return "✗ Libro no encontrado"
        
        if not libro.esta_disponible():
            return f"✗ No hay copias de {libro} disponibles"
        
        # Buscar usuario
        usuario = None
        for u in self._usuarios:
            if u.id_usuario == id_usuario:
                usuario = u
                break
        
        if not usuario:
            return "✗ Usuario no registrado"
\end{lstlisting}

\small(Continuación en siguiente slide...)
\end{frame}

\begin{frame}[fragile]{Proyecto: Biblioteca - Préstamo (cont.) \faBook}
\begin{lstlisting}
        if not usuario.puede_prestar():
            return f"✗ {usuario.nombre} alcanzó el límite de préstamos"
        
        # Realizar préstamo
        libro.copias_disponibles -= 1
        usuario.agregar_prestamo(libro)
        prestamo = Prestamo(libro, usuario)
        self._prestamos.append(prestamo)
        
        dias = (prestamo.fecha_devolucion - prestamo.fecha_prestamo).days
        return f"✓ Préstamo exitoso. Devolver antes de: {prestamo.fecha_devolucion.strftime('%d/%m/%Y')} ({dias} días)"
\end{lstlisting}
\end{frame}

\begin{frame}[fragile]{Proyecto: Biblioteca - Devolución \faUndoAlt}
\begin{lstlisting}
    def devolver_libro(self, titulo, id_usuario):
        # Buscar préstamo activo
        for prestamo in self._prestamos:
            if (prestamo.libro.titulo == titulo and 
                prestamo.usuario.id_usuario == id_usuario and 
                not prestamo.devuelto):
                
                # Procesar devolución
                prestamo.marcar_devuelto()
                prestamo.libro.copias_disponibles += 1
                prestamo.usuario.devolver_libro(prestamo.libro)
                
                vencido = " (¡VENCIDO!)" if prestamo.esta_vencido() else ""
                return f"✓ Libro devuelto{vencido}"
        
        return "✗ No se encontró préstamo activo"
\end{lstlisting}
\end{frame}

\begin{frame}[fragile]{Proyecto: Biblioteca - Reportes \faChartBar}
\begin{lstlisting}
    def listar_disponibles(self):
        print("\n=== LIBROS DISPONIBLES ===")
        for libro in self._catalogo:
            if libro.esta_disponible():
                print(f"• {libro} - {libro.copias_disponibles} copias")
    
    def prestamos_activos(self):
        print("\n=== PRÉSTAMOS ACTIVOS ===")
        for prestamo in self._prestamos:
            if not prestamo.devuelto:
                vencido = " ⚠ VENCIDO" if prestamo.esta_vencido() else ""
                print(f"• {prestamo}{vencido}")
    
    @staticmethod
    def validar_isbn(isbn):
        # Validación simple de ISBN-10
        return len(isbn) == 10 and isbn.isdigit()
\end{lstlisting}
\end{frame}

\begin{frame}[fragile]{Proyecto: Uso del sistema \faTerminal}
\begin{lstlisting}
# Crear biblioteca
bib = Biblioteca("Biblioteca Central")

# Agregar libros
libro1 = Libro("Cien años de soledad", "García Márquez", "1234567890", 2)
libro2 = Libro("Don Quijote", "Cervantes", "0987654321", 1)

bib.agregar_libro(libro1)
bib.agregar_libro(libro2)

# Registrar usuarios
user1 = Usuario("Ana García", "U001")
user2 = Usuario("Luis Pérez", "U002")

bib.registrar_usuario(user1)
bib.registrar_usuario(user2)
\end{lstlisting}

\small(Continuación en siguiente slide...)
\end{frame}

\begin{frame}[fragile]{Proyecto: Uso - Préstamos \faHandshake}
\begin{lstlisting}
# Realizar préstamos
print("\n" + "="*40)
print(bib.prestar_libro("Cien años", "U001"))
print(bib.prestar_libro("Don Quijote", "U002"))
print(bib.prestar_libro("Cien años", "U001"))

# Listar estado
bib.listar_disponibles()
bib.prestamos_activos()

# Devolver libro
print("\n" + "="*40)
print(bib.devolver_libro("Cien años de soledad", "U001"))

bib.listar_disponibles()
\end{lstlisting}
\end{frame}

\begin{frame}[fragile]{Proyecto: Salida \faArrowCircleRight}
\begin{tcolorbox}[colback=green!10,colframe=green,title={\small\faArrowCircleRight \ Salida},coltitle=fg0,left=3pt,right=3pt,top=2pt,bottom=2pt,boxsep=1pt]
\tiny\texttt{✓ Libro agregado: "Cien años de soledad" por García Márquez\\
✓ Libro agregado: "Don Quijote" por Cervantes\\
✓ Usuario registrado: Ana García (ID: U001)\\
✓ Usuario registrado: Luis Pérez (ID: U002)\\
\ \\
========================================\\
✓ Préstamo exitoso. Devolver antes de: 04/11/2025 (14 días)\\
✓ Préstamo exitoso. Devolver antes de: 04/11/2025 (14 días)\\
✗ No hay copias de "Cien años de soledad" por García Márquez disponibles\\
\ \\
=== LIBROS DISPONIBLES ===\\
• "Cien años de soledad" por García Márquez - 1 copias\\
\ \\
=== PRÉSTAMOS ACTIVOS ===\\
• "Cien años de soledad" por García Márquez - Ana García (ID: U001) [Activo]\\
• "Don Quijote" por Cervantes - Luis Pérez (ID: U002) [Activo]}
\end{tcolorbox}
\end{frame}

\begin{frame}{Proyecto alternativo: Sistema de inventario \faBoxes}
\begin{conceptbox}{Alternativa}
Sistema para gestionar inventario de productos en una tienda
\end{conceptbox}

\textbf{Clases sugeridas:}
\small
\begin{itemize}
    \item[\faBox] \textbf{Producto}: Base con nombre, precio, stock
    \item[\faTshirt] \textbf{ProductoFisico}: Hereda, agrega peso, dimensiones
    \item[\faLaptop] \textbf{ProductoDigital}: Hereda, agrega tamaño de archivo
    \item[\faShoppingCart] \textbf{Carrito}: Gestiona lista de compra
    \item[\faStore] \textbf{Tienda}: Catalogo, ventas, reportes
    \item[\faUser] \textbf{Cliente}: Datos, historial de compras
\end{itemize}
\end{frame}

\begin{frame}{Proyecto alternativo: Sistema de estudiantes \faGraduationCap}
\begin{conceptbox}{Otra alternativa}
Sistema para gestionar estudiantes, cursos y calificaciones
\end{conceptbox}

\textbf{Clases sugeridas:}
\small
\begin{itemize}
    \item[\faUser] \textbf{Persona}: Clase base con nombre, ID
    \item[\faUserGraduate] \textbf{Estudiante}: Hereda, agrega cursos, promedio
    \item[\faChalkboardTeacher] \textbf{Profesor}: Hereda, agrega cursos que imparte
    \item[\faBook] \textbf{Curso}: Nombre, código, créditos
    \item[\faClipboardList] \textbf{Calificacion}: Estudiante, curso, nota
    \item[\faUniversity] \textbf{Universidad}: Gestión completa
\end{itemize}
\end{frame}

\begin{frame}{Buenas prácticas POO \faCheckDouble}
\begin{conceptbox}{Principios SOLID}
\small
\begin{itemize}
    \item[\faCheck] \textbf{S}ingle Responsibility: Una clase, una responsabilidad
    \item[\faCheck] \textbf{O}pen/Closed: Abierto a extensión, cerrado a modificación
    \item[\faCheck] \textbf{L}iskov Substitution: Subclases sustituibles por su base
    \item[\faCheck] \textbf{I}nterface Segregation: Interfaces específicas
    \item[\faCheck] \textbf{D}ependency Inversion: Depender de abstracciones
\end{itemize}
\end{conceptbox}

\vspace{0.2cm}
\textbf{Otras prácticas:}
\small
\begin{itemize}
    \item[\faCheck] Nombres descriptivos para clases y métodos
    \item[\faCheck] Encapsular datos sensibles
    \item[\faCheck] Preferir composición sobre herencia cuando sea apropiado
    \item[\faCheck] Documentar clases con docstrings
\end{itemize}
\end{frame}

\begin{frame}[fragile]{Documentación con docstrings \faFileAlt}
\begin{lstlisting}
class CuentaBancaria:
    """
    Representa una cuenta bancaria básica.
    
    Attributes:
        titular (str): Nombre del titular de la cuenta
        saldo (float): Saldo actual de la cuenta
    
    Methods:
        depositar(monto): Deposita dinero en la cuenta
        retirar(monto): Retira dinero de la cuenta
    """
    
    def __init__(self, titular, saldo=0):
        """Inicializa una cuenta bancaria."""
        self.titular = titular
        self.saldo = saldo
\end{lstlisting}

\small(Continuación en siguiente slide...)
\end{frame}

\begin{frame}[fragile]{Documentación: Métodos \faBook}
\begin{lstlisting}
    def depositar(self, monto):
        """
        Deposita dinero en la cuenta.
        
        Args:
            monto (float): Cantidad a depositar
        
        Returns:
            str: Mensaje de confirmación
        
        Raises:
            ValueError: Si el monto es negativo
        """
        if monto < 0:
            raise ValueError("El monto no puede ser negativo")
        self.saldo += monto
        return f"Depósito exitoso. Saldo: ${self.saldo}"
\end{lstlisting}
\end{frame}

\begin{frame}{Resumen \faListUl}
\begin{conceptbox}{¿Qué aprendimos hoy?}
\small
\begin{itemize}
    \item[\faCheck] Conceptos de POO: clases y objetos
    \item[\faCheck] Constructor \_\_init\_\_ y atributos
    \item[\faCheck] Métodos de instancia, clase y estáticos
    \item[\faCheck] Encapsulación y convenciones de privacidad
    \item[\faCheck] Propiedades con @property
    \item[\faCheck] Métodos especiales: \_\_str\_\_, \_\_repr\_\_
    \item[\faCheck] Herencia y super()
    \item[\faCheck] Decoradores: @classmethod, @staticmethod
    \item[\faCheck] Proyecto: Sistema de gestión de biblioteca
\end{itemize}
\end{conceptbox}

\vspace{0.3cm}
\begin{center}
\large\textcolor{brown}{\faCube\ ¡POO organiza y potencia tu código!}
\end{center}
\end{frame}

\begin{frame}
\begin{center}
\Huge\textcolor{brown}{\faQuestionCircle}\\
\vspace{0.5cm}
\Huge ¡Preguntas!\\
\vspace{1cm}
\Large\textcolor{brownbright}{Gracias por su atención}
\end{center}
\end{frame}

\end{document}
